\section{The intermediate rows}\label{sec:intermediate}

Throughout this section abbreviate
\[
 e=\epsilon_0,\qquad k=k_0,\qquad l=l_0,\qquad
 j=j_0,\qquad p=p_0,\qquad \mu_{\rm HP}=\frac{16}{25},\qquad
 S_0=\frac{30869}{5000000}.
\]
The first five and $S_0$ are the fixed parameters from
\eqref{eq:old:2_5}; $\mu_{\rm HP}=16/25$ is set here for the 4P grid.

In this section the chart and all source predicates are those of
Section~\ref{sec:partition}.  If the selected chart is the $u$-chart, then
$(J,K;f,g;s)=(A,D;ac,bd;u)$.  If it is the $t$-chart, then
$(J,K;f,g;s)=(B,C;ad,bc;t)$.  In either case
\begin{equation}
 gK^2-fJ^2=s>0,\qquad fg=M,\qquad ut=L.
\AutoTag{eq:old:6_1}
\end{equation}
The chart choice is unique, including equality $u=t$.  Fix the residual
factors, the chart, the shift and the two selected corners.  Then the
inverse formulas in \eqref{eq:old:2_2} determine the original pair.  All
counts below keep the positivity, integrality, gcd, mask, shell and
first-match tests when a relaxed tuple is reconstructed.

\subsection{The supported-part row $4E$}

Let $f_J$ be the full prime-power part of $f$ supported on primes dividing
$J$.  Let $g_K$ be the full prime-power part of $g$ supported on primes
dividing $K$.  Write
\[
 \begin{aligned}
 f&=d_fv,\qquad g=d_gw,\qquad d_f=f_J,\qquad d_g=g_K,\\[-2pt]
 \operatorname{rad}(d_f)&\mid J,\qquad
 \operatorname{rad}(d_g)\mid K.
 \end{aligned}                                          \AutoTag{eq:old:6_2}
\]
Thus $d_f,d_g$ contain all powers of their supported primes.  For fixed
\((f,J)\) and \((g,K)\), each is uniquely determined.

Explicitly, \(d_f\) is the unique largest divisor of \(f\) all of whose
prime divisors divide \(J\), and \(v=f/d_f\).  Similarly, \(d_g\) is the
unique largest divisor of \(g\) supported on \(K\), and \(w=g/d_g\).  Once
\(g\) and \(K\) are recovered in the \(v\)-tag, \(d_g\) is determined by
\((g,K)\).  The same holds for \(d_f\) in the \(w\)-tag.  The mask assigns
every supported prime power to these determined parts.

In the \(u\)-chart, \(JK=AD=(n-m)/t\), so \(JK\le Yu/L\).  In the
\(t\)-chart, \(JK=BC=(m+n+1)/u\), so \(JK\le2Yt/L\).  We have
\(s=u\) in the first chart, \(s=t\) in the second, and \(ut=L\).  Hence
both inequalities give \(JK\le2Ys/L\).  The failed \(Q1\) and
selected-corner \(Q2\) tests give the weak lower bounds below.  Thus
\[
 JK\le \frac{2Ys}{L},\qquad
 L\ge Y^{1/3-l},\qquad J,K\ge Y^{1/3-j}.
\AutoTag{eq:old:6_3}
\]
Hence
\[
 J,K\le 2sY^{1/3+l+j}.                                  \AutoTag{eq:old:6_4}
\]
The row has $Y^{a_0}<s<Y^{b_0}$ and
\[
 \min(v,w)\le Y^p.
\AutoTag{eq:old:6_5}
\]
It has at most four disjoint tags.  They are the selected chart and either
$v\le Y^p$, or $v>Y^p$ and $w\le Y^p$.  By \eqref{eq:old:6_5}, the second
alternative covers the complement of the first.

Each tag uses a uniform supported-power bound.
For fixed $C,D,\alpha>0$, if $J\le Y^C$, then
\[
 \#\{d\le Y^D:\operatorname{rad}(d)\mid J\}
 \ll_{C,D,\alpha}Y^\alpha.                           \AutoTag{eq:old:6_6}
\]
To prove it, put $\sigma=\alpha/(2D)$.  Rankin's inequality gives
\[
 \#\{d\le Y^D:\operatorname{rad}(d)\mid J\}
 \le Y^{D\sigma}\prod_{q\mid J}(1-q^{-\sigma})^{-1}.
\]
Fix any $\delta>0$.  All primes above a fixed threshold satisfy
$-\log(1-q^{-\sigma})\le\delta\log q$.  The product over the finitely many
smaller primes is a constant.  The remaining product is at most
$J^\delta$.  Taking $\delta=\alpha/(2C)$ proves \eqref{eq:old:6_6},
uniformly in the moving label $J$.

In the tag $v\le Y^p$, enumerate $(s,v,J,d_f)$ and put
\[
 N=vd_fJ^2+s=fJ^2+s.
\]
For an actual source, $N=gK^2>0$.  Thus $K^2\mid N$, and each divisor
choice for $K^2$ forces $g=N/K^2$.  There are at most $\tauDiv(N)$ choices.
In the other tag enumerate $(s,w,K,d_g)$ and put
\[
 N=wd_gK^2-s=gK^2-s.
\]
Keep only \(N>0\).  On an actual source \(N=fJ^2\), so \(J^2\mid N\) and
$f=N/J^2$ is forced.  Discarding \(N\le0\) is safe, since an actual source
has \(N=fJ^2>0\).

Once $f,g$ are known, the ordered residual splittings are
\[
 (a,c)\text{ with }ac=f,\quad (b,d)\text{ with }bd=g
 \quad\text{or}\quad
 (a,d)\text{ with }ad=f,\quad (b,c)\text{ with }bc=g,
\]
according to the chart.  Their total number is at most $\tauDiv(f)\tauDiv(g)$.
Each splitting and $(s,J,K)$ determine the remaining corners, and then
$(m,n)$ by \eqref{eq:old:2_2}.  Candidates failing an original-source
predicate are discarded.  Once the ordered residuals are fixed, the
reconstruction has multiplicity at most one.  Then $ut=L$ determines the
opposite shift.

Before the divisor estimates, discard formal candidates that violate the
original source-height conditions.  Every retained source has
$J,K\le Y+1$, $f,g,d_f,d_g\le(Y+1)^2$, and
$1\le N\le(Y+1)^4$.  Apply \eqref{eq:old:6_6} to the one supported part.
Apply the standard estimate $\tauDiv(q_0)\ll_\rho q_0^\rho$ to each of $N,f,g$.
Fix a target loss $\eta>0$ and put $T=Y+1$. Apply \eqref{eq:old:6_6}
with $T$ in place of its displayed base variable $Y$, and take
$C=1$, $D=2$, and output exponent $\alpha=\eta/4$.  Its Rankin
parameters may be $\sigma=\eta/16$ and $\delta=\eta/8$.  For the divisor
bounds, take $\rho_N=\eta/16$ for $\tauDiv(N)$ and
$\rho_f=\rho_g=\eta/8$ for $\tauDiv(f),\tauDiv(g)$.  Since
$N\le T^4$ and $f,g\le T^2$, each of the four multiplicative factors is
$O_\eta(T^{\eta/4})$.  Their product is
$O_\eta(T^\eta)=O_\eta(Y^\eta)$, uniformly in the retained source
labels.  Thus each fixed $(s,v,J)$, or its counterpart $(s,w,K)$,
contributes at most $Y^\eta$ complete source records.  Sum over the small
endpoint, then over the selected corner using \eqref{eq:old:6_4}, then over
$s<Y^{b_0}$.  This gives
\[
 \#4E\ll_\eta
 Y^{p+1/3+l+j+\eta}\sum_{1\le s<Y^{b_0}}s
 \ll_\eta Y^{1/3+2b_0+l+j+p+\eta}.
\AutoTag{eq:old:6_7}
\]
The direct sum satisfies $\sum_{s<Y^{b_0}}s\ll Y^{2b_0}$.  At the parameter
values in \eqref{eq:old:2_5}
\[
 \frac13+2b_0+l+j+p=\frac{9907393}{15000000}=\frac23-S_0.
\AutoTag{eq:old:6_8}
\]
This equality, with $\eta$ arbitrary, is the claimed $4E$ bound.
\begin{rowproposition}[4E row estimate]\label{prop:row-4e}
For every \(\eta>0\), the complete first-match row satisfies
\[
 \#4E\ll_\eta Y^{2/3-S_0+\eta}.
\]
\end{rowproposition}
\begin{proof}
The full count is \eqref{eq:old:6_7}, and \eqref{eq:old:6_8} gives its exponent.
\end{proof}

\subsection{The opposite-shift cubic row $4H$}

For each prime dividing a residual, the radical mask places that prime in
one of the four corners.  The adjacent coprimalities determine which
corner.  Primes in \(a,b,c,d\) go respectively to one of
\((A,B),(C,D),(A,C),(B,D)\).  These choices are unique, since the two
corners in each listed pair are coprime.  Apply the same assignment to the
full prime powers.  This gives four pairwise coprime parts, assigned to
corners, whose product is \(M\):
\[
 (ac)_A,\quad (bd)_D,\quad (ad)_B,\quad (bc)_C.
\]
Set
\[
 G_u=\operatorname{rad}((ac)_A(bd)_D),\qquad
 G_t=\operatorname{rad}((ad)_B(bc)_C).
\]
Then \(G_uG_t=\operatorname{rad}(M)\).  In the \(u\)-chart, the supported
prime-power parts are \((ac)_A\) and \((bd)_D\).  Their radicals have
product \(G_u\).  In the \(t\)-chart the corresponding parts are
\((ad)_B\) and \((bc)_C\).  Their radicals have product \(G_t\).  Put
\[
 \Pi=\min(\operatorname{rad}(d_f),v_{\rm sel})
      \min(\operatorname{rad}(d_g),w_{\rm sel}).
\]
In \(4H\), \(v_{\rm sel},w_{\rm sel}>Y^p\) and \(\Pi\le Y^p\).  Hence
neither minimum can equal \(v_{\rm sel}\) or \(w_{\rm sel}\).  So
\(\Pi=\operatorname{rad}(d_f)\operatorname{rad}(d_g)\le Y^p\).  This is the
selected chart's assigned-corner radical product.  Since
\(\operatorname{rad}(M)>Y^{2/3-k}\), the opposite assigned-corner product
satisfies
\[
 G_{\rm opp}>R_0:=Y^{2/3-k-p}.                            \AutoTag{eq:old:6_9}
\]
The two-sided prime assignment in the mask makes \(G_{\rm opp}\) the
complementary assigned-corner product.

We use the following published theorem of Bombieri and Schmidt
\cite{BS1987}.  Let \(F\in\mathbb Z[X,Z]\) be an irreducible binary form
of degree \(r\ge3\), and let \(N\ge1\).  Then the number of primitive
integer pairs \((x,z)\) with \(|F(x,z)|=N\) is at most
\(C_{\rm BS} r^{1+\omega(N)}\), with an absolute \(C_{\rm BS}\).  Here
primitive means \(\gcd(x,z)=1\).  This is the unnumbered theorem on
pp.~69--70 of \emph{Inventiones Mathematicae} 88 (1987), 69--81, with the
correction on p.~445 of \emph{Inventiones Mathematicae} 97 (1989)
\cite{BS1989}. Its
bound is uniform in the coefficients.  It has no discriminant-coprimality
or coefficient-height hypothesis.

Hence, for each fixed \(C_0>0\) and every \(\eta>0\), uniformly for
positive integers \(c_1,c_2\le Y^{C_0}\) and
\(0<|h|\le Y^{C_0}\),
\[
 \#\{x,z>0:c_1x^3-c_2z^3=h\}\ll_{C_0,\eta}Y^\eta,
\AutoTag{eq:old:6_10}
\]
For the proof, let \(c=\gcd(c_1,c_2)\).  If \(c\nmid h\) there are no
solutions.  Otherwise divide by \(c\).  This gives coprime positive
coefficients \(c'_1,c'_2\) and \(h'=h/c\ne0\).  For each
\(g=\gcd(x,z)\), write \(x=gx_0,z=gz_0\).  Then \(g^3\mid h'\),
\((x_0,z_0)=1\), and
\(c'_1x_0^3-c'_2z_0^3=h'/g^3\).  If this cubic form is irreducible,
Bombieri--Schmidt gives at most
\(C3^{1+\omega(|h'|/g^3)}\) primitive pairs for that \(g\).  There are at
most \(\tauDiv(|h'|)\) choices of \(g\), and
\(3^{\omega(n)}\le\tauDiv(n)^{\log_2 3}\).  The divisor bound
\(\tauDiv(n)\ll_\rho n^\rho\), with \(\rho\) chosen after \(\eta\), makes
their total \(O_{C_0,\eta}(Y^\eta)\).

If the normalized cubic is reducible, \(c'_1/c'_2\) is a rational cube.
Coprimality then forces \(c'_1=A_{\rm cube}^3,c'_2=B_{\rm cube}^3\) for
coprime positive integers \(A_{\rm cube},B_{\rm cube}\).  Put
\(X=A_{\rm cube}x,Z=B_{\rm cube}z\).  Then
\[
 (X-Z)(X^2+XZ+Z^2)=h'.
\]
The nonzero integer $X-Z$ is a signed divisor of $h'$.  For each such
divisor the resulting quadratic has at most two integer roots.  The
divisor bound absorbs the content, gcd and reducible branches into
$Y^\eta$.  Here \(h'\ne0\), so the right side is nonzero.

For the opposite-chart conic count, write
\[
\begin{aligned}
 r_{\rm opp}&=\operatorname{rad}(d_{f,\rm opp}),&
 q_{\rm opp}&=\operatorname{rad}(d_{g,\rm opp}),\\
 d_{1,\rm opp}&=d_{f,\rm opp}/r_{\rm opp},&
 d_{2,\rm opp}&=d_{g,\rm opp}/q_{\rm opp},\\
 J_{\rm opp}&=r_{\rm opp}j_1,&K_{\rm opp}&=q_{\rm opp}z_1.
\end{aligned}
\]
The actual conic becomes the integral equation
\[
 (w_{\rm opp}d_{2,\rm opp}z_1^2)q_{\rm opp}^3
 -(v_{\rm opp}d_{1,\rm opp}j_1^2)r_{\rm opp}^3=s>0.
\AutoTag{eq:old:6_11}
\]
Here \(s\) is the shift of the opposite, counted chart. The selected
shift is then \(L/s\). The opposite chart has data
\[
(d_{f,\rm opp},d_{g,\rm opp},v_{\rm opp},w_{\rm opp})
\]
These are its supported parts and complementary factors from
\eqref{eq:old:2_4}. They differ from \((v_{\rm sel},w_{\rm sel})\) in
the \(4H\) predicate. Suppose the endpoint squares are at most
\(H_{4H}\) and \(r_{\rm opp}q_{\rm opp}\ge R_0\).  Fix
\(s,v_{\rm opp},w_{\rm opp},d_{1,\rm opp},d_{2,\rm opp},j_1,z_1\).
Apply \eqref{eq:old:6_10} to dyadic boxes
\(r_{\rm opp}\asymp R\), \(q_{\rm opp}\asymp Q\).  There are only
\(O(\log^2Y)\) such boxes, and \(RQ\ge R_0/4\).  The endpoint conditions
yield
\[
\begin{aligned}
 j_1&\le\sqrt{H_{4H}/(v_{\rm opp}d_{1,\rm opp}R^3)},\qquad
 z_1\le\sqrt{H_{4H}/(w_{\rm opp}d_{2,\rm opp}Q^3)},\\
 v_{\rm opp}w_{\rm opp}d_{1,\rm opp}d_{2,\rm opp}
 r_{\rm opp}q_{\rm opp}&<X_{\rm cap}.
\end{aligned}
\]
The displayed coefficient product is
\(v_{\rm opp}w_{\rm opp}d_{1,\rm opp}d_{2,\rm opp}r_{\rm opp}q_{\rm opp}=M\).
Failure of the \(O_R\) cut and \eqref{eq:old:2_7} give
\(M<\kappa_Y^2Y^{q_+}=\kappa_Y^2Y^{2/3+2e}\).  Define
\(X_{4H}:=\kappa_Y^2Y^{2/3+2e}\) here and take \(X_{\rm cap}=X_{4H}\).
If either upper bound is below one, that class is empty.  Otherwise put
\(N_{\rm opp}=v_{\rm opp}w_{\rm opp}\).  Sum over its at most
\(\tauDiv(N_{\rm opp})\) factorizations.  The number of \((j_1,z_1)\) is
at most the product of their upper bounds, so
\[
 \frac{H_{4H}}{\sqrt{d_{1,\rm opp}d_{2,\rm opp}}(RQ)^{3/2}}
 \sum_{N_{\rm opp}<X_{\rm cap}/(d_{1,\rm opp}d_{2,\rm opp}RQ)}
 \frac{\tauDiv(N_{\rm opp})}{\sqrt {N_{\rm opp}}}
 \ll \frac{H_{4H}\sqrt{X_{\rm cap}}\log(2X_{\rm cap})}
 {d_{1,\rm opp}d_{2,\rm opp}(RQ)^2}.
\]
Here we use $\sum_{N\le Z}\tauDiv(N)N^{-1/2}\ll\sqrt Z\log(2Z)$.  It
follows from $\tauDiv(N)=\sum_{ab=N}1$ by summing first over $b$.  Summing
\((d_{1,\rm opp},d_{2,\rm opp})\) harmonically and over the dyadic radical
boxes costs only \(Y^\eta\).  Since \(RQ\ge R_0/4\), the fixed-shift count
is
\[
 N_s\ll_\eta H_{4H}\sqrt{X_{4H}}R_0^{-2}Y^\eta.\AutoTag{eq:old:6_12}
\]
Fix an outer tuple
\(s,v_{\rm opp},w_{\rm opp},d_{1,\rm opp},d_{2,\rm opp},j_1,z_1\).
Then the cubic coefficients in \eqref{eq:old:6_11} are the fixed positive
integers \(c_1=w_{\rm opp}d_{2,\rm opp}z_1^2\) and
\(c_2=v_{\rm opp}d_{1,\rm opp}j_1^2\).
The cubic unknowns are \(q_{\rm opp}\) and \(r_{\rm opp}\).  So
\eqref{eq:old:6_10} applies uniformly in every dyadic radical box.
The coefficients may vary between outer tuples.  The uniformity of
\eqref{eq:old:6_10} permits the preceding sum over those tuples and boxes.
All canonical-source filters are kept on reconstruction.

In the following displays, \(J,K,s\) denote the counted chart and its
shift; this chart may be opposite to the first-match chart. The variables
in \eqref{eq:old:6_11} come from the canonical source. In the
$u$-chart they are
\[
 (f,g,J,K,s)=(ac,bd,A,D,u),
 \qquad (d_{f,\rm opp},d_{g,\rm opp},v_{\rm opp},w_{\rm opp})
 =(f_J,g_K,f/f_J,g/g_K),
\]
and in the $t$-chart they are
\[
 (f,g,J,K,s)=(ad,bc,B,C,t),
 \qquad (d_{f,\rm opp},d_{g,\rm opp},v_{\rm opp},w_{\rm opp})
 =(f_J,g_K,f/f_J,g/g_K).
\]
The two ordered residual factorizations are the only inverse choices.  In
the $u$-chart, fix a solution of \eqref{eq:old:6_11} and choose $ac=f$ and
$bd=g$.  Then
\[
 B=\frac{cA+bD}{u},\quad C=\frac{aA+dD}{u},\quad
 t=\frac{cd-ab}{u},\quad m=ABa,\quad n=ACc.                 \AutoTag{eq:old:6_12a}
\]
In the $t$-chart choose $ad=f$ and $bc=g$.  Then
\[
 A=\frac{dB-bC}{t},\quad D=\frac{cC-aB}{t},\quad
 u=\frac{cd-ab}{t},\quad m=ABa,\quad n=ACc.                 \AutoTag{eq:old:6_12b}
\]
Only positive integral reconstructions passing the original mask, gcd,
shell and first-match predicates are kept.  Hence
$\tauDiv(f)\tauDiv(g)$ bounds the inverse multiplicity, and $ut=L$
determines the other shift.

For the selected chart, $ut=L$ determines the opposite shift.
The source height and $R>Y^{5/3-e}$ give
\[
U=\kappa_Y Y^{1/3+e-a_0},\qquad
H_{4H}=12\kappa_Y^2Y^{4/3+e-a_0},\quad R_0=Y^{2/3-k-p},
\]
where \(\kappa_Y=1+Y^{-1}\le3/2\) for \(Y\ge2\).  If \(t<u\) is selected,
count the opposite $u$ chart with $H=uY$.  Then sum the actual $u<U$ once.
If $u\le t$ is selected, count the opposite $t$ chart with endpoint cap
$H_{4H}$.  Then sum the actual $t<U$ once; the other shift is forced.  For
the endpoint bound in the $t$-chart, the canonical identities give
\[
 (ad)B^2=\frac{m(n+1)}{ AD},\qquad
 (bc)C^2=\frac{n(m+1)}{ AD},\qquad
 (bc)C^2-(ad)B^2=t.                                      \AutoTag{eq:old:6_13a}
\]
The high-height condition \(ABCD>Y^{5/3-e}\) and
\(n(n+1)\le\kappa_Y Y^2\) give
\[
 cd=\frac{n(n+1)}{ ABCD}<\kappa_Y Y^{1/3+e}.
\]
Now \(BC=(m+n+1)/u\) and \(ABCD\,cd=n(n+1)\).  Also
\(n>Y/2\) and \(m+n+1\le3Y\).  Hence
\[
 AD=\frac{n(n+1)u}{ cd(m+n+1)}
   >\frac{uY^{2/3-e}}{12\kappa_Y}.                            \AutoTag{eq:old:6_13b}
\]
Thus both \(t\)-chart endpoint squares in \eqref{eq:old:6_13a} are at most
\(12\kappa_Y^2Y^{4/3+e}/u\).  On the subrow \(u>Y^{a_0}\) this is at most
\(H_{4H}\).  So this endpoint estimate uses both the high-height cut and
the subrow lower bound for \(u\).  The two counts are
\[
 \#4H_t\ll U^2Y\sqrt{X_{4H}}R_0^{-2}Y^\eta,\qquad
 \#4H_u\ll U H_{4H}\sqrt{X_{4H}}R_0^{-2}Y^\eta.
\AutoTag{eq:old:6_13}
\]
Here \(4H_t\) is the population selected in the \(t\)-chart
(\(t<u\)) and counted through the opposite \(u\)-chart.  Similarly,
\(4H_u\) is selected in the \(u\)-chart (\(u\le t\)) and counted through
the opposite \(t\)-chart.  The factor
\(\sqrt{X_{4H}}=\kappa_Y Y^{1/3+e}\) is the square-root factor in the
fixed-shift cubic count \eqref{eq:old:6_12}.  Substitute \(U\),
\(H_{4H}\) and \(R_0\) from above into either expression.  Both exponents
equal
\[
 \frac{2}{3}-2a_0+3e+2k+2p,
\]
so the saving is
\[
 r_H=2a_0-3e-2k-2p=\frac{140441}{5000000}>S_0.             \AutoTag{eq:old:6_14}
\]
The content split, reducible cubic count, residual divisor factors and
inverse formulas each take an arbitrary power loss.  Their number is fixed.
Summation over the shifts is the factor displayed in \eqref{eq:old:6_13}.
\begin{rowproposition}[4H row estimate]\label{prop:row-4h}
For every \(\eta>0\), the complete opposite-shift row satisfies
\[
 \#4H\ll_\eta Y^{2/3-140441/5000000+\eta}.
\]
\end{rowproposition}
\begin{proof}
The two endpoint counts in \eqref{eq:old:6_13} have the common exponent of
\eqref{eq:old:6_14}.  Their finite sum keeps that exponent.
\end{proof}

\subsection{The simultaneous-shift row $4P$}

The count in this row is simultaneous in both shifts.  Its source keeps
$v,w>Y^p$ and the strict defect bound $\Pi>Y^p$.  It also keeps the full
canonical source and all three high cuts.  It keeps the mask, unit and gcd
conditions and the selected-chart rule.  It keeps the outer and \(Q1/Q2\)
complements and the open middle-shift interval.  The 4P estimate below applies to this source.
Section~\ref{sec:long} counts the later $5W$ row separately.

In a dyadic box write
\[
 \begin{aligned}
 f&\asymp E,&g&\asymp F,&J&\asymp H_J,&K&\asymp H_K,\\
 S_h\le s<2S_h,&H_{4P}=2S_h,&X_{4P}&=H_J^2E,&V_*&=X_{4P}/H_K^2,\\
 N&=\lceil Y^{16/25}\rceil.
 \end{aligned}
\]
Here \(E\) and \(F\) are integer powers of two.  So \(EF\) is an integer
common denominator for the coefficient-lattice map below.  We also put
\[
 \begin{aligned}
 q&=\frac{2}{3}+2e,\qquad h_{\rm HP}=\frac{2}{3}+l+b_0,\\
 c&=1+e+l+b_0-2(1/3-j),\qquad
 \tau=\tau_0-\frac{1}{10000},\qquad \tau_0=1-\frac{3e}{2}.
 \end{aligned}                                          \AutoTag{eq:old:6_15}
\]
The same enlarged source gives $EF\ll Y^q$, $H_JH_K\ll Y^{h_{\rm HP}}$,
$E,F\ll Y^c$, and thickness $Y^{-\tau_0}$.  For
$z=J/K$, $w=g/f$, the conic is
\[
 gK^2-fJ^2=s.
\]
In particular \(H_{4P}/X_{4P}\ll Y^{-\tau_0}\), uniformly in the dyadic
box.  This is the thickness estimate used below.
For fixed nonnegative degrees \(A_{\rm det},B_{\rm det}\), use the
bidegree-\((A_{\rm det},B_{\rm det})\) columns
\[
 \{x^r y^j:0\le r\le A_{\rm det},\ 0\le j\le B_{\rm det}\},\qquad
 x=f/g,\quad y=K/J.                                      \AutoTag{eq:old:6_16}
\]
Clearing each row by \(g^{A_{\rm det}}J^{B_{\rm det}}\) gives the
integer entries in \eqref{eq:old:6_17c}.
The normalized determinant has the strict criterion
\[
\mu_{\rm HP}>\frac{9}{8}\frac{h_{\rm HP}q}{\tau},\qquad
\mu_{\rm HP}=\frac{16}{25}.
\AutoTag{eq:old:6_17}
\]
The statement of Lemma~\ref{lem:ratio-cell-determinant} fixes the 5W
parameters in \eqref{eq:old:7_2}--\eqref{eq:old:7_6}.  For 4P we use the
determinant calculation in its proof, with the separate 4P parameters.
The 4P grid exponent is \(a=\mu_{\rm HP}=16/25\), and
\[
 \frac35<a<\min(q_{\min},\tau_0),\qquad
 q_{\max}<\tau_0,\qquad b_0>\frac32e.
\]
The bounds \(H_J\le Y+1\) and \(E\le(Y+1)^2\) follow from the source
identities \eqref{eq:old:2_2}.  The 4P row itself has
\(H_{4P}>Y^{a_0}\).  The failed \(Q1\) and \(Q2\) cuts and the shell lower
bound \eqref{eq:old:2_8} supply the remaining endpoint restrictions in the
4P boxes.  In particular these source bounds give
\(X_{4P}/H_{4P}\le Y^4\) for sufficiently large \(Y\).  The Taylor order
is seven because \(8\mu_{\rm HP}=8(16/25)=128/25>4\).  With
\[
 p_1=\frac q{\mu_{\rm HP}}=\frac{5092607}{4800000}>1,
 \qquad
 p_2=\frac{h_{\rm HP}}\tau=\frac{24536965}{29719179},
\]
we have \(2/3<p_2<1\), \(p_1p_2<8/9\), and
\(p_1-1<3(1-p_2)^2\).  These include the strict trapezoid-branch hypotheses in
the proof of the general determinant lemma.  They give its 4P
specialization
\((a,p_1,p_2)=(16/25,5092607/4800000,24536965/29719179)\).
The determinant criterion needs these strict range and strip-weight
hypotheses as well as its displayed inequality.  The exact collared
determinant margin is positive; it is recorded just below.
Here \(\rho=p_2=h_{\rm HP}/\tau>2/3\), so 4P uses the trapezoid branch.
The stronger sufficient condition \(\mu_{\rm HP}>(9/8)q\rho\)
implies its criterion for \(2/3\le\rho<1\), because
\[
\frac98\rho\bigl(1+3(1-\rho)^2\bigr)-1
=\frac{(1-3(1-\rho))^3}{8}\ge0.
\]
At \eqref{eq:old:2_5}, the collared gap
\(\mu_{\rm HP}-(9/8)q\rho\) equals
$368942136449/39625572000000>0$.  Fix half this strict gap, then choose
the finite degree family.  All maximal-size integer minors have absolute
value less than one.  Hence they vanish, and the integerized matrix
\(\mathcal M\) has rank \(r<D\), where \(D\) is its number of columns.
If an interval contains fewer than \(D\) source rows, count those rows
directly; otherwise the vanishing \(D\)-minors imply rank below \(D\).
If \(r=0\), a coordinate vector is in its kernel.  Otherwise choose \(r\)
pivot columns and \(r\) rows giving a nonsingular minor \(A\), and one
free column \(j\).  Set the \(j\)-coordinate to \(\det A\), the other
free coordinates to zero, and the pivot coordinates to
\(-\operatorname{adj}(A)\mathcal M_{I,j}\).  This nonzero integer vector
annihilates the selected rows.  They span the row space, so it annihilates
every row.  Its entries are minors of bounded order and have polynomial
height.  The same argument covers the fixed anchor, the inverse monomials
and every endpoint branch.

For the normalization, use in each box the exact half-open ratio cells
\[
 r_i=\frac{i}{2N}\frac{H_J}{ H_K},\qquad
 I_i=(r_i,r_{i+1}],\qquad i=N,\ldots,4N-1 .                 \AutoTag{eq:old:6_17a}
\]
For a point in \(I_i\), put \(x=f/g\), \(y=K/J\), \(y_0=r_i^{-1}\),
\(\delta_{\rm cell}=s/(fJ^2)\), and \(\alpha_{\rm cell}=x/y_0^2-1\).
The conic gives the exact
identities
\[
 x=\frac{y^2}{1+\delta_{\rm cell}},\qquad
 \frac{y}{ y_0}=\sqrt{1+\alpha_{\rm cell}}\sqrt{1+\delta_{\rm cell}}.
\]
The cell width and \(s\le H_{4P}\) give
\[
 |\alpha_{\rm cell}|\le \frac{3}{ N},\qquad
 0\le \frac{y}{y_0}-\sqrt{1+\alpha_{\rm cell}}\le \frac{H_{4P}}{2X_{4P}}.       \AutoTag{eq:old:6_17b}
\]
Here \(r_{i+1}/r_i=1+1/i\le1+1/N\) and
\(1-(1+\delta_{\rm cell})^{-1/2}\le\delta_{\rm cell}/2\).  Since
\(X_2=X_{4P}/(2H_{4P})\), the conic error in \eqref{eq:old:6_17b} is
\(H_{4P}/(2X_{4P})=X_2^{-1}/4\).  Also
\(NH_{4P}/X_{4P}\le Y^{\mu_{\rm HP}-\tau_0+o(1)}<1\).  Choose a fixed
Taylor degree \(T\) and write
\(\sqrt{1+\alpha_{\rm cell}}=1+P_T(\alpha_{\rm cell})+\varepsilon_T\).
With \(X_1=N/3\), \eqref{eq:old:6_17b} gives
\(|\alpha_{\rm cell}|\le X_1^{-1}\). Choose \(T\)
large enough that the Taylor remainder is at most \(X_2^{-1}/2\), where
\(X_2=X_{4P}/(2H_{4P})\). Then the combined remainder
\(\beta=(y/y_0)-1-P_T(\alpha_{\rm cell})\) satisfies \(|\beta|\le X_2^{-1}\),
uniformly over all boxes and shifts.

Take columns \(x^r y^j\), \(0\le r\le A_{\rm det}\),
\(0\le j\le B_{\rm det}\), and multiply each row by
\(g^{A_{\rm det}}J^{B_{\rm det}}\).  The resulting entries are the integers
\[
 f^r g^{A_{\rm det}-r}K^jJ^{B_{\rm det}-j}.             \AutoTag{eq:old:6_17c}
\]
Write the polynomial in \eqref{eq:old:6_16} as
\[
C(x,y)=\sum_{r=0}^{A_{\rm det}}\sum_{j=0}^{B_{\rm det}}c_{rj}x^ry^j.
\]
With \(z=J/K=1/y\) and \(w=g/f=1/x\), its reciprocal is
\[
F(z,w)=z^{B_{\rm det}}w^{A_{\rm det}}C(1/w,1/z)
 =\sum_{r,j}c_{rj}z^{B_{\rm det}-j}w^{A_{\rm det}-r}.
\]
It has degrees at most \((B_{\rm det},A_{\rm det})\) in \((z,w)\) and the
same coefficient height.  Its coefficient vector permutes the full
rectangular coefficient array.  The column map is
\((r,j)\mapsto(a,b)=(B_{\rm det}-j,A_{\rm det}-r)\).  In particular,
\[
K^{B_{\rm det}}f^{A_{\rm det}}F(J/K,g/f)
 =J^{B_{\rm det}}g^{A_{\rm det}}C(f/g,K/J),
\]
so the two integer matrices agree after clearing denominators, up to that
column permutation.
After extracting the common anchor, the normalized entries are
\((1+\alpha_{\rm cell})^r(1+P_T(\alpha_{\rm cell})+\beta)^j\).  If
\(D_c=(A_{\rm det}+1)(B_{\rm det}+1)\),
the anchor and row-clearing factors have logarithmic cost at most
\[
 \frac{D_c}{2}\{A_{\rm det}\Lambda_1+B_{\rm det}\Lambda_2\},\qquad
 \Lambda_1=\log(324EV_*),\quad\Lambda_2=\log(8H_JH_K),
\AutoTag{eq:old:6_17d}
\]
The column-anchor factor is \(y_0^{2r+j}\).  Its exponent sum is
\[
\sum_{\substack{0\le r\le A_{\rm det}\\0\le j\le B_{\rm det}}}(2r+j)
=D_c\left(A_{\rm det}+\frac{B_{\rm det}}2\right).
\]
The other anchor bound is
\[
J y_0^{1/2}<2H_J(2H_K/H_J)^{1/2}=(8H_JH_K)^{1/2}.
\]
Also \(V_*=E H_J^2/H_K^2\), so
\[
18V_*\frac{H_K}{H_J}=18\sqrt{EV_*},\qquad
(18\sqrt{EV_*})^2=324EV_*.
\]
These inequalities give the constants \(8\) and \(324\) in the uniform
anchor bounds.
The conic and the thickness give \(EV_*\asymp EF\).  So replacing one by
the other changes only an absolute logarithmic constant.  The displayed
cost follows from \(g<9V_*\), \(y_0\le2H_K/H_J\), and \(J<2H_J\).
It is independent of \(i\) and \(s\).

The proof of Lemma~\ref{lem:ratio-cell-determinant} derives the finite
Cauchy--Binet expansion, the strip-weight lower bound, the branch
integrals and the integer-minor conclusion.  It does so for general
logarithmic parameters satisfying its strict inequalities.  We apply that
calculation with the 4P parameters and effective scales below.
For this row put
\[
 X_1=N/3,\qquad X_2^{\rm eff}=Y^\tau,
\]
and choose fixed constants \(C_1,C_2\) such that, uniformly in every
box,
\[
 \Lambda_1^{\rm eff}=q\log Y+C_1\ge\log(324EV_*),\qquad
 \Lambda_2^{\rm eff}=h_{\rm HP}\log Y+C_2\ge\log(8H_JH_K).
\]
Such constants exist by the source bounds \(EV_*\ll EF\ll Y^q\) and
\(H_JH_K\ll Y^{h_{\rm HP}}\).  Also \(H_{4P}/X_{4P}\ll Y^{-\tau_0}\), where
\(\tau_0>\tau\).  So for all sufficiently large \(Y\),
\(X_2=X_{4P}/(2H_{4P})\ge X_2^{\rm eff}\).  Thus the actual Taylor error
\(|\beta|\le X_2^{-1}\) is at most \((X_2^{\rm eff})^{-1}\).  The
arithmetic cost in \eqref{eq:old:6_17d} is bounded by the one obtained from
\(\Lambda_1^{\rm eff},\Lambda_2^{\rm eff}\).  With
\[
 \ell_1=\log X_1,\quad \ell_2^{\rm eff}=\tau\log Y,\quad
 P_{\rm eff}=\frac{\Lambda_1^{\rm eff}}{\ell_1},\quad
 R_{\rm eff}=\frac{\Lambda_2^{\rm eff}}{\ell_2^{\rm eff}},\quad
 \ell_{\rm det}=\frac{\ell_1}{\ell_2^{\rm eff}},
\]
we have
\[
 P_{\rm eff}=\frac q{\mu_{\rm HP}}+O(1/\log Y),\quad
 R_{\rm eff}=\frac{h_{\rm HP}}\tau+O(1/\log Y),\quad
 \ell_{\rm det}=\frac{\mu_{\rm HP}}\tau+O(1/\log Y).
\]
At the limiting 4P vector,
\[
 \mu_{\rm HP}=\frac{16}{25},\qquad
 p_1=\frac{5092607}{4800000},\qquad
 p_2=\frac{24536965}{29719179},\qquad
 \ell_{\rm det}=\frac{6400000}{9906393}.
\]
The corresponding normalized determinant margins are
\[
 \frac89-p_1p_2=\frac{368942136449}{28530411840000}>0,\qquad
 3(1-p_2)^2-(p_1-1)
 =\frac{14253084474184329457}{471055786898155200000}>0.
\]
Hence the strict margin in \eqref{eq:old:6_17}, with the fixed collar
\(\tau=\tau_0-1/10000\), applies to these uniform majorants
for all sufficiently large \(Y\).  The finite determinant calculation uses
only the displayed upper bounds for the two errors and the arithmetic
cost.  So it gives one degree choice and one polynomial-height bound,
uniformly in the cell, shift and box.

After bihomogenization and factorization, we count each factor.
A factor that is not absolutely irreducible meets a distinct conjugate
in only boundedly many rational projective points at fixed degree.  For
corner-degree-zero factors, first fix the primitive coefficient pair,
then its ordered residual factors.  These fix $L=cd-ab$.  Discard $L\le0$
and enumerate the divisor choices $s\mid L$.  For each fixed
$(f,g,s)$ apply Theorem~\ref{thm:fibre} with
$(a,b,h,n,m)=(f,g,s,J,K)$, $x=2Y$, and fixed height exponent
$A_{\rm ht}=2$.
The residuals are at most $Y+1$, whence $f,g,s\le(Y+1)^2$ and
$J,K\le Y+1$.  This costs $Y^\eta$ per coefficient label.  The selected
shift is enumerated among the divisors of \(L\).  Explicitly, this cost
consists of three parts.  They are the ordered residual splitting factors
\(\tauDiv(f)\tauDiv(g)\), the selected-shift choices \(\tauDiv(L)\), and
the uniform square-fibre bound for each fixed \((f,g,s)\).  Give each
operation its own small enough internal exponent.  Then their product is
\(Y^\eta\).

For a coefficient-degree-zero factor, fix its primitive positive corner
pair $(J,K)$.  For two source coefficient pairs the conic gives
\[
 K^2(g_1f_2-g_2f_1)=s_1f_2-s_2f_1,
 \qquad |g_1f_2-g_2f_1|
 \le \frac{4H_{4P}E}{H_K^2}
 \ll EF\frac{H_{4P}}{X_{4P}}
 \ll Y^{q-\tau_0}=o(1).
\AutoTag{eq:old:6_18}
\]
Here $s_i<H_{4P}$, $f_i<2E$, $K\ge H_K$, $X_{4P}=H_J^2E$, and
 $V_*=X_{4P}/H_K^2\ll F$ by the source box bounds.  The cross determinant
is an integer.  It vanishes for all sufficiently large $Y$, since
$\tau_0-q=9351751/30000000>0$.  Primitivity of $(f_i,g_i)$ then forces
the coefficient pairs to agree, and the conic fixes the shift.  This
branch contributes \(O(N)\) primitive conic-coordinate tuples across the
primary cells.  To return to complete canonical source records, include
the ordered residual factorizations.  They cost
\(\tauDiv(f)\tauDiv(g)\ll_\eta Y^\eta\), so this branch is
\(O_\eta(Y^\eta N)\).

Let the coefficient degree be $b_q\ge1$ and the corner degree $a_c\ge2$.
Then apply Heath--Brown's bihomogeneous Lemma~2~\cite{HB2012} to the
factor after the unimodular coefficient-lattice change of basis.  Its
bidegree is $(b_q,a_c)$.  The first pair is the primitive transformed
coefficient pair of radius $O(T_1+1)$.  The second pair is the original
primitive corner pair $(J,K)$, with no size hypothesis.  The factor height
is $Y^C$ for a fixed $C$.  Thus the count is
$O((T_1+1)^{2/a_c+\varepsilon}Y^{C\varepsilon})$.  Here $2/a_c\le1$ and
$T_1+1\le Y^{C'}$ for a fixed $C'$.  Choose the theorem's internal
$\varepsilon<\eta/(C+C'+1)$; this bounds the count by $Y^\eta(T_1+1)$.
For corner degree one write the factor as
$KP(f,g)+JQ(f,g)$ with coprime homogeneous $P,Q$.  Their nonzero resultant
and homogeneous Bezout identities give
$d_0=\gcd(P(f,g),Q(f,g))\mid\operatorname{Res}(P,Q)$, since
$\gcd(f,g)=1$.  The factor equation forces
$(J,K)=\pm(P(f,g)/d_0,-Q(f,g)/d_0)$, with at most one positive sign.
Write $f_J=d_f$, $g_K=d_g$ for the supported parts and
$v=f/d_f$, $w=g/d_g$.
Write \(\ell_f,\ell_g\) for the endpoint labels in the four patterns
\((P(0,1),Q(1,0))\):
\[
\begin{array}{c|c|c|c}
P(0,1)&Q(1,0)&\ell_f&\ell_g\\ \hline
\ne0&\ne0&\rad(d_f)&\rad(d_g)\\
\ne0&=0&\rad(d_f)&w\\
=0&\ne0&v&\rad(d_g)\\
=0&=0&v&w .
\end{array}
\AutoTag{eq:old:6_18a}
\]
If $P(0,1)\ne0$, reducing $d_0J=P(f,g)$ modulo $\rad(d_f)$ gives
$\rad(d_f)\mid P(0,1)$.  If $P(0,1)=0$, then $f\mid P(f,g)$, and
$\gcd(v,J)=1$ gives $v\mid d_0\mid\operatorname{Res}(P,Q)$.  The same
argument with $Q(1,0)$ proves the two $g$-side alternatives.  The chosen
endpoint labels satisfy \(\ell_f\ell_g\ge\Pi>Y^p\), since each is at
least the corresponding minimum in the definition of \(\Pi\).  Each label
divides a fixed nonzero endpoint coefficient or resultant of polynomial
height.  So there are \(\ll_\eta Y^\eta\) possible pairs of endpoint
labels, for any prescribed $\eta>0$.  Also \(\gcd(\ell_f,\ell_g)=1\),
since \(\gcd(f,g)=1\) and the labels are supported on the respective parts
of \(f\) and \(g\).  Take two distinct primitive coefficient pairs with
the same labels.  Their divisibilities combine, and the nonzero cross
determinant is divisible by \(\ell_f\ell_g\).  Take a coefficient cell of
width $O(F/(EN))$ in $g/f$.  There the ratios are separated by at least
$\ell_f\ell_g/(4E^2)$, since $f_i<2E$.  Thus the number of pairs in one
cell is
\[
 O\left(1+\frac{EF}{ N \ell_f\ell_g}\right).
\]
Summing over the $O(N)$ coefficient cells and the divisor labels gives
\[
 \ll_\eta Y^\eta\left(N+\frac{EF}{ \ell_f\ell_g}\right)
 \ll_\eta Y^\eta\left(N+\frac{EF}{ Y^p}\right),
\]
including unit labels and endpoint equalities. A factor equal to a source
coordinate vanishes only when that coordinate is zero, so it has no points
in the positive source. This is separate from the lattice shortest-vector
axes, which are retained.
This branch counts the primitive coefficient and corner tuples.
Recovering complete source pairs still needs the ordered residual
splittings of \(f\) and \(g\), at most \(\tauDiv(f)\tauDiv(g)\).  That
multiplicity is absorbed in the displayed \(Y^\eta\).

For the high-degree factors use a second, coefficient-ratio grid.  Write
\[
 r_{i_c}=\frac FE\frac{i_c}{N},\qquad
 C_{i_c}=(r_{i_c},r_{i_c+1}],\qquad
 \lfloor N/2\rfloor\le i_c<2N.
\AutoTag{eq:old:6_19a}
\]
The dyadic source has \(g/f>F/(2E)\), since \(g\ge F\) and \(f<2E\).
Thus these \(O(N)\) cells cover its full range
\(F/(2E)<g/f<2F/E\).  Now use
\[
 \frac{g}{ f}=\left(\frac{J}{ K}\right)^2+\frac{s}{ fK^2},
 \qquad \frac{s}{gK^2}\ll \frac{H_{4P}}{X_{4P}}\ll Y^{-\tau_0},
\AutoTag{eq:old:6_19b}
\]
and $\tau_0>\mu_{\rm HP}$.  The relative displacement from the square map is
then $o(1/N)$.  Each grid has relative width $O(1/N)$, and squaring
preserves relative widths on positive intervals.  Hence each coefficient
cell meets $O(1)$ primary $J/K$ cells, and conversely.  Thus the cell
polynomials from \eqref{eq:old:6_17a} give $O(1)$ factors per coefficient
cell.

For a coefficient cell define the rational linear map
\[
 \mathcal T_{i_c}(g,f)=
 \left(\frac{N}{ F}(g-r_{i_c}f),\frac{f}{ E}\right).
\AutoTag{eq:old:6_19c}
\]
For a point in this cell, $f/E\in[1,2)$ and
\[
0\le \frac{N}{ F}(g-r_{i_c}f)
 <\frac{N}{ F}\,2E\,\frac{F}{ EN}=2,
\]
so its image lies in $[0,2]\times[1,2]$.  The matrix of
$\mathcal T_{i_c}$ is
\[
\begin{pmatrix}N/F&-i_c/E\\0&1/E\end{pmatrix},\qquad
\det\mathcal T_{i_c}=N/(EF).
\]
Let $\Lambda_{i_c}=\mathcal T_{i_c}\mathbb Z^2$.  Choose a
Lagrange-reduced basis $v_1,v_2$ of this image lattice.  Its lengths are
$\lambda_1\le\lambda_2$, and its angle is bounded away from zero.  Write
$v_j=\mathcal T_{i_c}h_j$.  Since $v_1,v_2$ are a lattice basis,
$h_1,h_2\in\mathbb Z^2$ form a unimodular basis of $\mathbb Z^2$.
Hence every primitive coefficient pair has a primitive transformed
pair $(x_1,x_2)$ in $(g,f)=x_1h_1+x_2h_2$.  Put
$T_1=1/\lambda_1$, $T_2=1/\lambda_2$.  The determinant and reduction give
\[
 T_1\ge T_2,\qquad T_1T_2\asymp \frac{EF}{ N},\qquad
 T_1\gg\sqrt{EF/N},\qquad
 \max(|x_1|,|x_2|)\ll T_1
\AutoTag{eq:old:6_19d}
\]
for the transformed coordinates $(x_1,x_2)$.  For the last bound, the
bounded image rectangle and the angle bound imply
$|x_j|\ll1/\lambda_j\le T_1$.  The source has
$EF\gg Y^{2/3-e}$ and $N\asymp Y^{16/25}$, so $EF/N$ grows as a positive
power.  A common denominator for the matrix of $\mathcal T_{i_c}$ is
$EF$.  So every nonzero image lattice vector has length at least
$(EF)^{-1}$.  Also $\lambda_1\ll\sqrt{N/(EF)}$ by the two-dimensional
Minkowski bound, and $\lambda_1\lambda_2\asymp N/(EF)$ by reduction.
Hence $\lambda_2\ll N$.  The inverse matrix is
\[
\mathcal T_{i_c}^{-1}(y_1,y_2)
=\bigl((F/N)(y_1+i_cy_2),Ey_2\bigr).
\]
Since $i_c\asymp N$ and $E,F,N\le Y^C$ for a fixed $C$, both $v_j$
and their integer preimages $h_j$ have polynomial height.  So substitution
by the unimodular matrix $[h_1\ h_2]$ keeps every fixed-degree factor at
polynomial coefficient height.  The unimodular change uses no moving
denominator.

The census below includes axes and first points.  At scale
$D_{\rm lat}\le T_1<2D_{\rm lat}$, let $(g_{\rm lat},f_{\rm lat})$ be the
primitive preimage of the shortest basis vector.  From
\eqref{eq:old:6_19c},
\[
 |f_{\rm lat}|\ll E/D_{\rm lat},\quad |g_{\rm lat}|\ll F/D_{\rm lat},\quad
 |i_cf_{\rm lat}-(NE/F)g_{\rm lat}|\ll E/D_{\rm lat}.
\AutoTag{eq:old:6_19e}
\]
If $f_{\rm lat},g_{\rm lat}\ne0$, then $E/D_{\rm lat}$ and $F/D_{\rm lat}$
are at least a fixed positive constant.  There are $O(F/D_{\rm lat})$
choices for $g_{\rm lat}$.  For each, the integer product $i_cf_{\rm lat}$
lies in an interval of length $O(E/D_{\rm lat})$.  This gives that many
possible values.  Each nonzero value has at most
$\tauDiv(|i_cf_{\rm lat}|)\ll_\eta Y^\eta$ factorizations into
$i_c f_{\rm lat}$, since all factors have polynomial height.  Thus
\[
 \#\{i_c:T_1\asymp D_{\rm lat},\ f_{\rm lat}g_{\rm lat}\ne0\}
 \ll_\eta Y^\eta \frac{EF}{D_{\rm lat}^2}.
\AutoTag{eq:old:6_19f}
\]
If $f_{\rm lat}=0$, primitivity gives $g_{\rm lat}=\pm1$ and
$T_1\asymp F/N$.  So the sum over $O(N)$ cells is $O(F+N)$.  If
$g_{\rm lat}=0$, then $f_{\rm lat}=\pm1$, $i_c\asymp N$, and
$T_1\asymp E/N$, giving $O(E+N)$.  Cells with $T_1<1$ contribute $O(N)$.
These cells, both axes and every $+1$ are kept.  Weight
\eqref{eq:old:6_19f} by $T_1\asymp D_{\rm lat}$ and sum geometrically
from $D_{\rm lat}\gg\sqrt{EF/N}$.  This yields
\[
 \sum_{i_c}(T_1(i_c)+1)\ll_\eta
 Y^\eta\{N+\sqrt{EFN}+E+F\}.
\AutoTag{eq:old:6_19}
\]
This count factors the exact integer product \(i_cf_{\rm lat}\). The number
of cells in this dyadic range is
\(O_\eta(Y^\eta EF/D_{\rm lat}^2)\).

For a factor with $b_q\ge1$ and $a_c\ge2$, the Heath--Brown bound above
is $Y^\eta(T_1+1)$.  Summing over coefficient cells by \eqref{eq:old:6_19}
gives $Y^\eta(N+\sqrt{EFN}+E+F)$.  Short matrices, nonabsolute factors and
the two zero-degree branches contribute $Y^\eta N$.  The corner-linear
branch contributes $Y^\eta(N+EF/Y^p)$, by the endpoint-label spacing
argument above.  These cases exhaust the fixed-degree factorization.

Combine the lattice census \eqref{eq:old:6_19} with the corner-linear
count.  The complete bound then has the form
\[
 \#4P\ll_\eta Y^\eta\left(N+\sqrt{EFN}+E+F+\frac{EF}{Y^p}\right).
\]
These local bounds are summed over the \(O(\log^4Y)\) choices of dyadic
scales and the \(O(\log Y)\) shift scales.  The resulting logarithmic
factor is absorbed into the arbitrary \(Y^\eta\) loss.
The four distinct exponent contributions, in the order
\(EF/Y^p\), \(\sqrt{EFN}\), \(N\), and \(E+F\), are
\[
 q-p=\frac{9907393}{15000000},\quad
 \frac{q+\mu}{2}=\frac{9892607}{15000000},\quad
 \mu=\frac{16}{25},\quad
 c=\frac{5277821}{10000000}.
\AutoTag{eq:old:6_20}
\]
Their savings relative to $2/3$ are respectively
\[
S_0,\qquad \frac{107393}{15000000},\qquad \frac{2}{75},\qquad
 \frac{4166537}{30000000}.
\AutoTag{eq:old:6_21}
\]
The minimum is $S_0=30869/5000000$.  We combine the uniform lattice census
\eqref{eq:old:6_19}, the finite factor family and the polynomial-height
estimates.  With the divisor losses and the one simultaneous source
inversion, this gives
\[
\#4P\ll_\eta Y^{2/3-S_0+\eta}.
\AutoTag{eq:old:6_22}
\]
\begin{rowproposition}[4P row estimate]\label{prop:row-4p}
For every \(\eta>0\), the complete simultaneous-shift row satisfies
\[
 \#4P\ll_\eta Y^{2/3-S_0+\eta}.
\]
\end{rowproposition}
\begin{proof}
Equations \eqref{eq:old:6_19}--\eqref{eq:old:6_21} bound every factor
branch and its exponent.  Then \eqref{eq:old:6_22} sums them with the
source-uniform lattice and inverse bounds.
\end{proof}

The $4P$ estimate applies to its enlarged source.  Section~\ref{sec:long}
counts $5W$ from its original-cell source and geometric envelope.
